\chapter{Emission}\label{chap:emission}

The \texttt{emit-fixed} pipeline applies transforms T1 (action narrowing),
T3 (resource narrowing) to each Allow statement, threading through a
need-set.  The composition theorem guarantees the emitted policy
provably never grants what the input did not.

\section{Definitions}

\begin{definition}[emitFixed]\label{def:emitFixed}
\lean{emitFixed}
\leanok
Given a policy and a need-set, applies \texttt{transformStmt}
(narrowAction then narrowResource) to each statement via
\texttt{filterMap}.  Deny statements pass through unchanged.
\end{definition}

\section{Pipeline Theorem}

\begin{theorem}[emitFixed\_narrows]\label{thm:emitFixed_narrows}
\uses{thm:filterMap_narrows, thm:stmtGrantsAction_narrow, thm:matchPattern_cs_literal_eq}
\lean{emitFixed_narrows}
\leanok
\textbf{Pipeline theorem (T1+T3).}
If \texttt{allows (emitFixed p ns).1 req ctx = true} and every need
action is wildcard-free, then \texttt{allows p req ctx = true}.
The emitted policy never grants access the original did not.

The proof factors through \texttt{filterMap\_narrows}, showing that
\texttt{transformStmt} preserves Deny statements unchanged and
produces match-monotone, effect/condition-preserving replacements for
Allow statements.  Action monotonicity uses
\texttt{stmtGrantsAction\_narrow}; resource monotonicity uses
\texttt{matchPattern\_cs\_literal\_eq}.
\end{theorem}

\section{Completeness}

\begin{corollary}[grants\_complete]\label{cor:grants_complete}
\uses{thm:grants_complete_generic}
\lean{grants_complete}
\leanok
If an IAM policy allows a request, there exists an Allow statement
that matches and whose condition is not~F.  Direct application of
the generic \texttt{grants\_complete} via \texttt{iamSem}.
\end{corollary}
